\chapter{Le transformeur de vision, l'œil de la machine}
\label{chap:vit}

\begin{quote}\itshape
Nous savons ce qu'est une image pour un ordinateur, nous savons ce que veut
dire apprendre, et nous savons jouer au jeu des trous. Il nous manque une
chose, et une seule : l'appareil qui regarde. Dans ce chapitre, nous le
construisons pièce par pièce, sans en sauter aucune. À la fin, vous saurez
exactement ce qui se passe entre le moment où une vidéo entre dans la
machine et le moment où elle en ressort sous forme de résumé.
\end{quote}

\begin{prerequis}
Les chapitres~\ref{chap:image-video}, \ref{chap:apprendre}
et~\ref{chap:auto-supervision}. Il faut savoir qu'une image est une grille
de nombres et connaître sa forme ; il faut savoir qu'un modèle possède des
boutons que l'on tourne pour réduire une erreur ; et il faut se souvenir
qu'un encodeur transforme une image en un petit paquet de nombres appelé
représentation. Nous n'aurons besoin de rien d'autre. Les seules opérations
employées seront l'addition, la multiplication et la division.
\end{prerequis}

Au chapitre précédent, nous avons dit qu'une machine cache un morceau de ce
qu'elle regarde et essaie de le deviner. Nous avons même écrit un petit
programme qui le faisait. Mais nous avons triché sur un point : notre monde
tenait dans une grille de huit cases sur huit, et notre modèle était une
simple couche de calcul. Une vraie vidéo, elle, contient des millions de
nombres. Il faut donc une machine à regarder beaucoup plus solide.

Cette machine porte un nom : le \emph{transformeur de vision}. Tout le
monde l'appelle par son surnom anglais, \emph{ViT}, qui se prononce
«~vi-té~» et veut dire \emph{Vision Transformer}. C'est l'œil de V-JEPA~2.1,
et c'est aussi l'œil de presque tous les grands modèles d'images
d'aujourd'hui. Nous allons le monter devant vous, morceau par morceau.

\section{Pourquoi on ne peut pas regarder les pixels un par un}

\begin{enclair}
L'idée du transformeur, c'est de laisser chaque morceau de l'image regarder
tous les autres morceaux. Le problème, c'est que si les morceaux sont des
pixels, il y en a beaucoup trop : le nombre de regards à échanger devient
si grand qu'aucune machine au monde ne pourrait le calculer.
\end{enclair}

Commençons par la question qui commande tout le reste. Pour comprendre une
image, il ne suffit pas de regarder chaque endroit tout seul dans son coin.
Il faut les mettre en rapport les uns avec les autres.

Prenons un exemple. Sur une photo, vous voyez une oreille pointue en haut à
gauche, des moustaches au milieu, et une queue touffue en bas à droite.
Aucun de ces trois morceaux ne dit «~chat~» à lui tout seul : une oreille
pointue, cela peut être un renard ; des moustaches, cela peut être un phoque.
C'est le \emph{rapprochement} des trois qui donne la réponse. Une bonne
machine à regarder doit donc permettre à chaque endroit de l'image de
consulter tous les autres endroits.

Comptons maintenant le prix de cette consultation. Imaginons que chaque
pixel consulte chaque autre pixel. Dans une image de $256 \times 256$, il y
a $65\,536$ pixels. Le nombre de consultations vaut alors

\[
  65\,536 \times 65\,536 = 4\,294\,967\,296 ,
\]

soit plus de quatre milliards de consultations, et cela pour \emph{une
seule} image, à \emph{un seul} endroit du modèle. Un clip de seize images
contient seize fois plus de pixels, donc $16 \times 16 = 256$ fois plus de
consultations encore. C'est tout simplement impossible.

\begin{intuition}
Imaginez une salle où chaque personne doit regarder toutes les autres, une
par une. À dix personnes, cela fait $10 \times 10 = 100$ regards : on s'en
sort en une minute. À mille personnes, cela fait $1\,000 \times 1\,000 =
1\,000\,000$ de regards : la réunion durerait des semaines. Le nombre de
personnes a été multiplié par cent, mais le travail a été multiplié par dix
mille. C'est ce qu'on appelle un coût \emph{qui grimpe au carré}, et c'est
l'ennemi numéro un de tout ce chapitre.
\end{intuition}

\begin{definition}[Un coût qui grimpe au carré][cout-carre]
On dit qu'un coût \emph{grimpe au carré} quand il suit le nombre d'éléments
multiplié par lui-même. Appelons $n$ le nombre d'éléments — écrire une
lettre à la place d'un nombre permet simplement de parler de tous les cas à
la fois. Si chacun de ces $n$ éléments doit consulter tous les autres, le
nombre de consultations vaut $n \times n$, ce qui s'écrit aussi $n^2$ et se
lit «~$n$ au carré~».
\end{definition}

\begin{exemple}[Voir le coût s'emballer]
Mettons quelques nombres côte à côte, et regardons la colonne de droite
grandir :

\begin{center}
\footnotesize
\begin{tabular}{rr}
\toprule
\textbf{Nombre d'éléments} & \textbf{Travail à faire} \\
\midrule
$10$    & $100$ \\
$100$   & $10\,000$ \\
$1\,000$ & $1\,000\,000$ \\
\bottomrule
\end{tabular}
\end{center}

À chaque fois que la colonne de gauche est multipliée par dix, celle de
droite est multipliée par cent. Voilà pourquoi on dit que le coût
\emph{s'emballe} : il grandit bien plus vite que ce qu'on lui donne.
\end{exemple}

La solution est celle que tout le monde emploie quand une réunion est trop
nombreuse : au lieu de faire parler chaque personne, on fait parler chaque
\emph{table}. On regroupe. C'est exactement ce que fait le transformeur de
vision, et c'est le sujet de la section suivante.

\section{Le patch : découper l'image en petits carrés}

\begin{enclair}
Plutôt que de laisser dialoguer les pixels un par un, on découpe l'image en
petits carrés de taille fixe. Chaque carré devient un porte-parole, et ce
sont ces porte-parole qui dialoguent entre eux. Il y en a beaucoup moins, et
tout devient calculable.
\end{enclair}

\begin{definition}[Le patch][patch]
Un \emph{patch} est un petit carré de pixels découpé dans l'image, toujours
de la même taille. Le mot vient de l'anglais et veut dire «~pièce~» ou
«~morceau~», comme une pièce de tissu cousue sur un pantalon. Dans
V-JEPA~2.1, un patch mesure $16 \times 16$ pixels. On dit aussi que la
\emph{taille de patch} vaut $16$.
\end{definition}

Le découpage se fait sans laisser de trou et sans que deux morceaux se
chevauchent : on pose une grille régulière sur l'image, exactement comme on
découpe une plaque de chocolat en carrés le long des rainures. Chaque carré
est ensuite mis à la queue leu leu, dans un ordre convenu une fois pour
toutes : la première ligne de gauche à droite, puis la deuxième ligne, et
ainsi de suite. On obtient une longue file de morceaux.

\begin{schema}{Le découpage en patchs, puis la mise en file. À gauche,
l'image est recouverte d'une grille régulière ; chaque case est un patch et
porte son numéro. À droite, les mêmes patchs sont rangés les uns derrière
les autres, dans l'ordre de lecture d'un texte : ligne par ligne, de gauche
à droite. Cette file est ce que la machine manipulera désormais.}[decoupage-patchs]
  \begin{scope}[x=1cm, y=1cm]
    \node[font=\sffamily\small\bfseries, text=mbSlate] at (2.0,4.5)
         {L'image et sa grille};
    \foreach \ri in {0,...,3}{
      \foreach \ci in {0,...,3}{
        \pgfmathtruncatemacro{\num}{\ri*4+\ci+1}
        \pgfmathtruncatemacro{\tone}{14+4*\num}
        \fill[mbAzure!\tone] (\ci*1.0,3.9-\ri*1.0) rectangle ++(1.0,-1.0);
        \draw[draw=mbSlate!55, line width=0.6pt]
              (\ci*1.0,3.9-\ri*1.0) rectangle ++(1.0,-1.0);
        % au-delà de la moitié de la grille, le fond devient trop foncé
        % pour du texte sombre : on bascule en blanc.
        \ifnum\num>10\relax
          \node[font=\sffamily\scriptsize\bfseries, text=white]
                at (\ci*1.0+0.5,3.9-\ri*1.0-0.5) {\num};
        \else
          \node[font=\sffamily\scriptsize\bfseries, text=mbInk]
                at (\ci*1.0+0.5,3.9-\ri*1.0-0.5) {\num};
        \fi
      }
    }
    \node[font=\sffamily\scriptsize, text=mbSlate] at (2.0,-0.45)
         {$4 \times 4 = 16$ patchs};

    \draw[fleche accent] (4.35,1.9) -- (5.55,1.9);
    \node[etiquette, above] at (4.95,1.98) {on met en file};

    \begin{scope}[xshift=5.9cm, yshift=1.55cm]
      \node[font=\sffamily\small\bfseries, text=mbSlate] at (2.56,2.95)
           {La file de patchs};
      \foreach \k in {0,...,4}{
        \pgfmathtruncatemacro{\num}{\k+1}
        \pgfmathtruncatemacro{\tone}{14+4*\num}
        \fill[mbAzure!\tone] (\k*0.72,0) rectangle ++(0.72,0.72);
        \draw[draw=mbSlate!55, line width=0.6pt]
              (\k*0.72,0) rectangle ++(0.72,0.72);
        \node[font=\sffamily\scriptsize\bfseries, text=mbInk]
              at (\k*0.72+0.36,0.36) {\num};
      }
      \node[font=\sffamily\small, text=mbSlate] at (4.05,0.36) {$\cdots$};
      \fill[mbAzure!78] (4.4,0) rectangle ++(0.72,0.72);
      \draw[draw=mbSlate!55, line width=0.6pt] (4.4,0) rectangle ++(0.72,0.72);
      \node[font=\sffamily\scriptsize\bfseries, text=white] at (4.76,0.36) {16};
      \draw[accolade bas] (0,-0.12) -- (5.12,-0.12);
      \node[font=\sffamily\scriptsize, text=mbSlate] at (2.56,-0.75)
           {une file de $16$ morceaux};
    \end{scope}
  \end{scope}
\end{schema}

Combien de morceaux obtient-on ? Le calcul est le plus simple du monde : on
divise la largeur par la taille du patch, on divise la hauteur par la taille
du patch, et on multiplie les deux résultats.

\begin{exemple}[Compter les patchs d'une image de V-JEPA 2.1]
Pendant la phase principale de l'entraînement, les images font
$256 \times 256$ pixels et les patchs mesurent $16 \times 16$. Sur la
largeur, on peut donc poser

\[
  \frac{256}{16} = 16 \text{ patchs,}
\]

et autant sur la hauteur. Le nombre total de patchs vaut

\[
  16 \times 16 = 256 .
\]

Retenez bien ce nombre, $256$ : nous allons le retrouver sans cesse.
Souvenez-vous aussi que cette image contenait $65\,536$ pixels. Nous venons
donc de ramener le nombre d'éléments à consulter de $65\,536$ à $256$,
c'est-à-dire de le diviser par $256$. Et comme le coût grimpe au carré, le
travail, lui, est divisé par $256 \times 256$, soit par plus de soixante-cinq
mille.
\end{exemple}

\begin{piege}
La division doit tomber juste. Si vous prenez une image de $250$ pixels de
côté avec des patchs de $16$, la division $250/16$ ne donne pas un nombre
entier, et il vous restera une bande de pixels sur le bord que personne ne
sait où ranger. C'est pour cela que les tailles d'images choisies dans
l'article sont toujours des multiples de $16$ : $256$, $384$, $512$. Quand
vous changerez de résolution dans vos propres essais, vérifiez toujours ce
point avant de lancer quoi que ce soit.
\end{piege}

\begin{enia}
Ce découpage est la toute première opération de V-JEPA~2.1, celle qui a lieu
avant tout le reste. L'article lui donne un nom que nous retrouverons tout
au long de la partie~IV : le \emph{tokeniseur}. C'est simplement la pièce
qui fabrique les morceaux.
\end{enia}

Nous avons donc une file de petits carrés. Mais un petit carré, c'est encore
une grille de nombres. Or un transformeur ne sait pas manipuler des grilles :
il manipule des listes de nombres. Il faut donc transformer chaque carré en
une liste. C'est ce que nous faisons tout de suite.

\section{Du carré au jeton : le plongement}

\begin{enclair}
Chaque petit carré est d'abord mis à plat, c'est-à-dire écrit comme une
simple liste de nombres. Puis on lui applique un calcul qui le remplace par
une autre liste, plus courte et mieux rangée. Cette nouvelle liste
s'appelle un jeton, et c'est elle que le modèle manipulera du début à la
fin.
\end{enclair}

Reprenons un patch de $16 \times 16$ pixels en couleur. Combien de nombres
contient-il ? Nous savons le calculer depuis le chapitre~\ref{chap:image-video} :
la hauteur, fois la largeur, fois les trois canaux de couleur.

\[
  16 \times 16 \times 3 = 768 \text{ nombres.}
\]

On écrit ces $768$ nombres les uns derrière les autres, en une seule ligne.
Cette opération s'appelle \emph{mettre à plat}, ou \emph{aplatir}. Rien n'est
perdu, rien n'est ajouté : on change seulement la façon de ranger.

Ensuite vient le vrai travail. On ne garde pas ces $768$ nombres bruts : on
les remplace par une autre liste, obtenue grâce à des \emph{coefficients}.
Un coefficient, c'est tout simplement un nombre par lequel on en multiplie
un autre. Ceux dont nous parlons ici sont \emph{appris} : ce sont exactement
les boutons du chapitre~\ref{chap:apprendre}. Au début ils valent n'importe
quoi, et l'entraînement les règle petit à petit.

\begin{definition}[Le plongement, ou embedding][embedding]
Un \emph{plongement} est une liste de nombres qui représente un morceau de
donnée. On dit aussi \emph{embedding}, qui est le mot anglais employé
partout. On l'obtient en multipliant les nombres du morceau par des
coefficients appris, puis en additionnant les résultats. La longueur de cette
liste s'appelle la \emph{dimension} du plongement : c'est un choix que fait
celui qui construit le modèle, et elle est la même pour tous les morceaux.
\end{definition}

\begin{schema}{Du carré au jeton, en deux gestes. On part d'un patch de
$16 \times 16$ pixels en couleur. On le met d'abord à plat, c'est-à-dire
qu'on écrit ses nombres les uns derrière les autres : cela en fait $768$, et
aucun n'a été touché. On leur applique ensuite un calcul dont les
coefficients sont appris, et qui rend une liste bien plus courte. Cette
liste est le jeton.}[carre-vers-jeton]
  \begin{scope}[x=1cm, y=1cm]
    %---- 1. le patch, centré sur la ligne y = 0 ---------------------
    \foreach \ri in {0,...,3}{
      \foreach \ci in {0,...,3}{
        \pgfmathtruncatemacro{\tone}{18+8*(\ri+\ci)}
        \fill[mbAzure!\tone] (\ci*0.45,0.9-\ri*0.45) rectangle ++(0.45,-0.45);
        \draw[draw=mbSlate!45, line width=0.4pt]
              (\ci*0.45,0.9-\ri*0.45) rectangle ++(0.45,-0.45);
      }
    }
    \node[font=\sffamily\scriptsize, text=mbSlate, align=center] at (0.9,-1.3)
         {un patch\\ $16 \times 16 \times 3$};

    %---- 2. on met à plat -------------------------------------------
    \draw[fleche accent] (2.1,0) -- (3.3,0);
    \node[etiquette, above] at (2.7,0.08) {on met à plat};

    \foreach \k in {0,...,8}{
      \pgfmathtruncatemacro{\tone}{18+9*\k}
      \fill[mbAzure!\tone] (3.6+\k*0.32,-0.22) rectangle ++(0.32,0.44);
      \draw[draw=mbSlate!45, line width=0.4pt]
            (3.6+\k*0.32,-0.22) rectangle ++(0.32,0.44);
    }
    \node[font=\sffamily\small, text=mbSlate] at (6.72,0) {$\cdots$};
    \node[font=\sffamily\scriptsize, text=mbSlate, align=center] at (5.4,-1.3)
         {$768$ nombres,\\ aucun n'a changé};

    %---- 3. le calcul appris ----------------------------------------
    \draw[fleche accent] (7.1,0) -- (8.3,0);
    \node[etiquette, above] at (7.7,0.08) {calcul appris};

    \foreach \k in {0,...,3}{
      \pgfmathtruncatemacro{\tone}{30+16*\k}
      \fill[mbTeal!\tone] (8.6+\k*0.32,-0.22) rectangle ++(0.32,0.44);
      \draw[draw=mbTeal!80!black, line width=0.5pt]
            (8.6+\k*0.32,-0.22) rectangle ++(0.32,0.44);
    }
    \node[font=\sffamily\scriptsize, text=mbTeal!55!black, align=center]
         at (9.24,-1.3) {le jeton :\\ des mesures,\\ pas des pixels};

    %---- l'accolade qui montre que la liste a raccourci --------------
    \draw[accolade] (3.6,0.32) -- (7.04,0.32);
    \draw[accolade] (8.6,0.32) -- (9.88,0.32);
    \node[annotation] at (9.24,0.95) {plus court};
  \end{scope}
\end{schema}

\begin{definition}[Le jeton, ou token][jeton]
Un \emph{jeton} est le plongement d'un patch, c'est-à-dire la liste de
nombres qui remplace ce patch dans toute la suite des calculs. On dit aussi
\emph{token}, en anglais. Le mot est bien choisi : comme un jeton de casino
remplace une somme d'argent, un jeton remplace ici un morceau d'image. Le
découpage suivi du plongement s'appelle la \emph{tokenisation}.
\end{definition}

\begin{exemple}[Un plongement en tout petit]
Pour voir le calcul en entier, prenons un patch minuscule de $2 \times 2$
pixels en niveaux de gris, donc quatre nombres, et une dimension de
plongement de $2$ seulement. Le patch vaut

\[
  \begin{pmatrix} 10 & 20 \\ 30 & 40 \end{pmatrix}
  \quad \longrightarrow \quad
  \text{mis à plat : } (10,\ 20,\ 30,\ 40).
\]

Le modèle possède, pour cette opération, huit coefficients appris : quatre
pour fabriquer le premier nombre du jeton, quatre pour le second. Disons
qu'ils valent

\[
  (0{,}1 \; ; \; 0{,}0 \; ; \; -0{,}1 \; ; \; 0{,}2)
  \quad \text{et} \quad
  (0{,}0 \; ; \; 0{,}5 \; ; \; 0{,}0 \; ; \; -0{,}5).
\]

Le premier nombre du jeton s'obtient en multipliant terme à terme puis en
additionnant :

\[
  10 \times 0{,}1 + 20 \times 0{,}0 + 30 \times (-0{,}1) + 40 \times 0{,}2
  = 1 + 0 - 3 + 8 = 6 .
\]

Le second nombre, de la même façon :

\[
  10 \times 0{,}0 + 20 \times 0{,}5 + 30 \times 0{,}0 + 40 \times (-0{,}5)
  = 0 + 10 + 0 - 20 = -10 .
\]

Le jeton vaut donc $(6\ ;\ -10)$. Quatre nombres sont devenus deux nombres,
et ces deux nombres ne sont plus des intensités de gris : ce sont des
mesures fabriquées par le modèle. Regardez la seconde liste de
coefficients, $(0 \; ; \; 0{,}5 \; ; \; 0 \; ; \; -0{,}5)$ : elle ajoute
le pixel en haut à droite et retranche celui d'en bas à droite. Elle mesure
donc une différence entre le haut et le bas. Personne ne le lui a demandé :
c'est l'entraînement qui a réglé ces quatre nombres ainsi.
\end{exemple}

\begin{intuition}
Le plongement est un résumé. Là où les $768$ nombres bruts disaient «~voici
l'éclat de chacun de mes pixels~», le jeton dit plutôt «~je contiens un bord
oblique, plutôt sombre, plutôt vert~». Personne n'a écrit ces mots nulle
part : ce sont les coefficients appris qui décident, tout seuls, de ce
qu'il vaut la peine de mesurer.
\end{intuition}

\begin{surlesdonnees}
Dans V-JEPA~2.1, ce calcul n'est pas fait patch par patch dans une boucle,
mais d'un seul coup par une opération appelée \emph{convolution}. Une
convolution, c'est une petite fenêtre que l'on fait glisser sur l'image :
à chaque arrêt, elle regarde ce qui se trouve dessous et en tire quelques
nombres. Ici, la fenêtre mesure $16 \times 16$ pixels et elle avance de $16$
pixels à chaque pas. Avancer d'exactement la largeur de la fenêtre, c'est
précisément découper sans chevauchement. Le résultat est le même que ce que
nous venons de calculer à la main, mais il est obtenu bien plus vite. Nous
reverrons la convolution de près au chapitre~13 ; pour l'instant, retenez
seulement l'image de la fenêtre qui glisse.
\end{surlesdonnees}

Nous savons transformer une image en une file de jetons. Reste une question
que nous avons soigneusement mise de côté : et le temps ?

\section{Le tubelet : découper aussi dans le temps}

\begin{enclair}
Une vidéo, c'est une pile d'images. Si l'on découpait chaque image
séparément, on obtiendrait beaucoup trop de jetons. On découpe donc aussi
dans l'épaisseur de la pile : chaque morceau prend deux images d'un coup.
Ce petit pavé s'appelle un tubelet.
\end{enclair}

\begin{definition}[Le tubelet][tubelet]
Un \emph{tubelet} est un patch qui possède en plus une épaisseur dans le
temps. Au lieu d'être un carré plat découpé dans une seule image, c'est un
petit pavé qui traverse plusieurs images de suite. Dans V-JEPA~2.1, ce pavé
mesure $16 \times 16$ pixels et $2$ images d'épaisseur : on dit que la
\emph{taille de tubelet} vaut $2$.
\end{definition}

\begin{schema}{Le patch et le tubelet. À gauche, un patch : un carré plat,
découpé dans une seule image. À droite, un tubelet : le même carré, pris sur
deux images qui se suivent. Les traits pointillés montrent qu'il s'agit bien
du même endroit. Le tubelet contient donc deux fois plus de pixels, et il
porte en lui un petit bout de mouvement, puisqu'il voit cet endroit à deux
instants.}[patch-tubelet]
  \begin{scope}[x=1cm, y=1cm]
    %================ à gauche : le patch, plat =======================
    \node[font=\sffamily\small\bfseries, text=mbSlate] at (1.2,3.65)
         {Un patch : plat};
    \foreach \ri in {0,...,3}{
      \foreach \ci in {0,...,3}{
        \draw[draw=mbSlate!40, fill=white, line width=0.4pt]
              (\ci*0.6,2.4-\ri*0.6) rectangle ++(0.6,-0.6);
      }
    }
    \fill[mbBlue!45] (0.6,1.8) rectangle ++(0.6,-0.6);
    \draw[draw=mbBlue, line width=1pt] (0.6,1.8) rectangle ++(0.6,-0.6);
    \node[font=\sffamily\scriptsize, text=mbSlate] at (1.2,-0.45)
         {une seule image};

    % le trait qui sépare les deux moitiés du dessin
    \draw[draw=mbSlate!25, line width=0.5pt] (3.6,-0.9) -- (3.6,3.35);

    %================ à droite : le tubelet ===========================
    % Les deux images sont posées côte à côte, sans se recouvrir : ainsi
    % la case retenue reste visible sur l'une comme sur l'autre.
    \node[font=\sffamily\small\bfseries, text=mbSlate] at (7.5,3.65)
         {Un tubelet : le même carré, sur deux images};

    \foreach \img/\dx/\teinte in {1/4.6/55, 2/8.0/35}{
      \begin{scope}[shift={(\dx,0)}]
        \foreach \ri in {0,...,3}{
          \foreach \ci in {0,...,3}{
            \draw[draw=mbSlate!40, fill=white, line width=0.4pt]
                  (\ci*0.6,2.4-\ri*0.6) rectangle ++(0.6,-0.6);
          }
        }
        \fill[mbTeal!\teinte] (0.6,1.8) rectangle ++(0.6,-0.6);
        \draw[draw=mbTeal!85!black, line width=1pt]
              (0.6,1.8) rectangle ++(0.6,-0.6);
        \node[font=\sffamily\scriptsize, text=mbSlate] at (1.2,-0.45)
             {image \img};
      \end{scope}
    }

    % les pointillés : c'est bien le même endroit des deux côtés
    \draw[draw=mbTeal!85!black, line width=0.7pt, dashed]
          (5.8,1.8) -- (8.6,1.8);
    \draw[draw=mbTeal!85!black, line width=0.7pt, dashed]
          (5.8,1.2) -- (8.6,1.2);

    \draw[accolade bas] (4.6,-0.75) -- (10.4,-0.75);
    \node[font=\sffamily\scriptsize, text=mbTeal!55!black] at (7.5,-1.4)
         {le tubelet, c'est les deux carrés ensemble};

    \draw[-{Stealth[round,length=2.2mm,width=1.7mm]},
          draw=mbAmber!85!black, line width=1pt] (4.6,2.72) -- (10.4,2.72);
    \node[font=\sffamily\scriptsize, text=mbAmber!80!black, above]
         at (7.5,2.75) {le temps passe};
  \end{scope}
\end{schema}

Pourquoi coller deux images ensemble plutôt que de les traiter séparément ?
Pour deux raisons, et les deux sont bonnes.

La première est une simple affaire de nombre. Un clip de seize images
découpé image par image donnerait $16 \times 256 = 4\,096$ jetons. Avec des
tubelets de deux images, on en obtient deux fois moins. Comme le coût grimpe
au carré, le travail est divisé par quatre.

La seconde raison est plus intéressante. Un tubelet voit le même endroit à
deux instants différents. Il contient donc, dès sa fabrication, une petite
trace de ce qui a bougé à cet endroit. Le mouvement n'est pas quelque chose
que le modèle devra reconstituer plus tard : il est déjà là, dans le
morceau.

\begin{exemple}[Compter les jetons d'un clip de la phase principale]
Pendant la phase principale, un clip contient $16$ images de
$256 \times 256$ pixels. Comptons dans les trois directions.

Dans le temps, avec des tubelets de $2$ images :
\[ \frac{16}{2} = 8 \text{ tranches.} \]

En hauteur, avec des patchs de $16$ pixels :
\[ \frac{256}{16} = 16 \text{ rangées.} \]

En largeur, de même : $16$ colonnes. Le nombre total de jetons vaut donc le
produit des trois :

\[
  8 \times 16 \times 16 = 2\,048 \text{ jetons.}
\]

Voilà le deuxième nombre à retenir par cœur : \textbf{$2\,048$ jetons pour
un clip de la phase principale}.
\end{exemple}

\begin{exemple}[Compter les jetons d'un clip de refroidissement]
Vous connaissez déjà cette phase. Le chapitre~\ref{chap:image-video} vous a
donné la taille de ses clips, $64$ images de $384 \times 384$ pixels, et le
chapitre~\ref{chap:apprendre} vous a dit d'où lui vient son nom : sur les
douze mille dernières étapes, le taux d'apprentissage redescend jusqu'à
presque zéro, et le modèle se pose en douceur. Refaisons pour elle le calcul
des jetons :

\[
  \frac{64}{2} = 32, \qquad \frac{384}{16} = 24, \qquad \frac{384}{16} = 24,
\]
\[
  32 \times 24 \times 24 = 18\,432 \text{ jetons.}
\]

C'est le troisième nombre à retenir. Comparons-le au précédent :

\[
  \frac{18\,432}{2\,048} = 9 .
\]

Neuf fois plus de jetons : c'est exactement le rapport que nous avions
trouvé au chapitre~\ref{chap:image-video} en comptant les pixels bruts, et
ce n'est pas un hasard, puisque le découpage garde les proportions.

Mais voici ce que le chapitre~\ref{chap:image-video} ne pouvait pas encore
nous dire. Là-bas, neuf fois plus de nombres voulait dire neuf fois plus de
mémoire, et nous étions passés de $1{,}6$ à $14{,}5$ gigaoctets par lot. Ici,
c'est le \emph{calcul} qui est en jeu, et le calcul grimpe au carré :
\[ 9 \times 9 = 81 . \]
Un clip de refroidissement coûte donc \emph{quatre-vingt-une fois} plus de
travail qu'un clip de la phase principale. La mémoire était multipliée par
neuf, le calcul l'est par quatre-vingt-un. Voilà pourquoi l'article ne fait
durer cette phase que $12\,000$ pas, alors que la phase principale en compte
$135\,000$.
\end{exemple}

\begin{aretenir}
Quatre lignes à garder en tête pour tout le livre :

\begin{center}
\footnotesize
\begin{tabular}{llr}
\toprule
\textbf{Entrée} & \textbf{Calcul} & \textbf{Jetons} \\
\midrule
Image $256 \times 256$        & $16 \times 16$               & $256$ \\
Image $512 \times 512$        & $32 \times 32$               & $1\,024$ \\
Clip $16 \times 256 \times 256$ & $8 \times 16 \times 16$    & $2\,048$ \\
Clip $64 \times 384 \times 384$ & $32 \times 24 \times 24$   & $18\,432$ \\
\bottomrule
\end{tabular}
\end{center}

La règle est toujours la même : on divise chaque côté par la taille du
découpage, puis on multiplie les résultats entre eux.
\end{aretenir}

Nos morceaux sont prêts et rangés en file. Mais en les mettant en file, nous
avons perdu quelque chose d'important, et il faut le récupérer sans tarder.

\section{Dire où l'on se trouve : l'encodage de position}

\begin{enclair}
Une fois les jetons mis en file, plus rien ne dit lequel venait du coin en
haut à gauche et lequel venait du milieu. On doit donc écrire l'adresse de
chaque morceau à l'intérieur même de son jeton.
\end{enclair}

Voici le problème, et il est sérieux. Le calcul de la section suivante,
l'attention, traite la file de jetons comme un \emph{sac} : il ne regarde
jamais l'ordre. Mélangez la file, il rendra exactement le même résultat. Or
une image dont on mélange les morceaux n'est plus la même image du tout.

\begin{intuition}
Prenez un puzzle terminé, photographiez-le, puis jetez toutes les pièces
dans une boîte. Les pièces sont toujours là, aucune n'a été perdue, et
pourtant le dessin a disparu. Ce qui manque, ce n'est pas la matière, c'est
la \emph{place} de chaque pièce. Voilà exactement ce que nous devons rendre
au modèle.
\end{intuition}

\begin{definition}[L'encodage de position][encodage-position]
L'\emph{encodage de position} est la façon dont on inscrit, dans le jeton
lui-même, l'endroit d'où venait son morceau. Sans lui, le modèle ne saurait
pas distinguer le haut du bas, la gauche de la droite, ni le début de la fin
d'une vidéo.
\end{definition}

V-JEPA~2.1 emploie une méthode particulière, appelée \emph{RoPE}. Ce nom est
une abréviation anglaise de \emph{Rotary Position Embedding}, ce qui veut
dire «~encodage de position par rotation~». Le mot rotation dit tout : au
lieu d'ajouter un numéro au jeton, on \emph{tourne} le jeton d'un angle qui
dépend de sa position.

\begin{schema}{L'idée de RoPE, expliquée avec une aiguille d'horloge. On
prend les nombres du jeton deux par deux et on les voit comme une flèche.
Cette flèche est ensuite tournée d'un angle proportionnel à la position du
morceau : rien à la position $0$, un cran à la position $1$, deux crans à la
position $2$, et ainsi de suite. La longueur de la flèche ne change jamais :
seule sa direction porte l'information de place.}[rope-horloge]
  \begin{scope}[x=1cm, y=1cm]
    \foreach \k/\ang in {0/0, 1/30, 2/60, 3/90}{
      \begin{scope}[shift={(\k*3.4,0)}]
        \draw[draw=mbSlate!30, line width=0.5pt] (0,0) circle (1.05);
        \draw[draw=mbSlate!25, line width=0.4pt] (-1.25,0) -- (1.25,0);
        \draw[draw=mbSlate!25, line width=0.4pt] (0,-1.25) -- (0,1.25);
        % la position de départ, rappelée en pointillé sur chaque cadran
        \draw[draw=mbSlate!45, line width=0.6pt, dashed] (0,0) -- (1.05,0);
        % l'angle parcouru depuis le départ
        \ifnum\ang>0\relax
          \draw[draw=mbAmber!70, line width=0.7pt] (0.45,0) arc (0:\ang:0.45);
        \fi
        \draw[vecteur accent] (0,0) -- (\ang:1.05);
        \fill[mbInk] (0,0) circle (1.4pt);
        \node[font=\sffamily\scriptsize\bfseries, text=mbSlate] at (0,1.55)
             {position \k};
        \node[font=\sffamily\scriptsize, text=mbAmber!80!black] at (0,-1.62)
             {tourné de $\ang^\circ$};
      \end{scope}
    }
  \end{scope}
\end{schema}

\begin{exemple}[Tourner une flèche, avec les nombres]
Prenons deux nombres d'un jeton, disons $(1\ ;\ 0)$, et voyons-les comme une
flèche qui pointe vers la droite. Convenons que chaque position fait tourner
la flèche de trente degrés. Mesurons à la règle où arrive la pointe de la
flèche après chaque rotation. Voici ce que l'on relève :

\begin{center}
\footnotesize
\begin{tabular}{crr}
\toprule
\textbf{Position} & \textbf{Angle} & \textbf{Les deux nombres} \\
\midrule
$0$ & $0^\circ$  & $(1{,}000 \ ; \ 0{,}000)$ \\
$1$ & $30^\circ$ & $(0{,}866 \ ; \ 0{,}500)$ \\
$2$ & $60^\circ$ & $(0{,}500 \ ; \ 0{,}866)$ \\
$3$ & $90^\circ$ & $(0{,}000 \ ; \ 1{,}000)$ \\
\bottomrule
\end{tabular}
\end{center}

Les deux nombres changent, alors que le contenu du morceau, lui, n'a pas
bougé. C'est la position, et elle seule, qui a fait tourner la flèche.
\end{exemple}

Pourquoi tourner plutôt qu'ajouter un numéro ? Parce que la rotation a un
avantage qui tombe à pic. Quand deux jetons se comparent, ce qui compte
n'est pas leur angle à chacun, mais l'\emph{écart} entre leurs deux angles.
Or cet écart ne dépend que de la \emph{distance} entre leurs positions.

\begin{exemple}[Ce qui compte, c'est l'écart]
Reprenons nos trente degrés par position. Deux jetons placés aux positions
$1$ et $3$ sont séparés de deux crans, donc de $60$ degrés. Deux autres
jetons placés aux positions $4$ et $6$ sont eux aussi séparés de deux crans,
donc des mêmes $60$ degrés. Pour le modèle, les deux paires sont dans la
même situation.

C'est très exactement ce que l'on veut. Un bord vertical reste un bord
vertical, qu'il se trouve à gauche ou à droite de la photo. Ce qui importe,
c'est la place des choses \emph{les unes par rapport aux autres}, et non
l'endroit exact où elles se trouvent dans l'image.
\end{exemple}

Dans une vidéo, il y a trois directions à repérer : la ligne, la colonne et
l'instant. On fait donc tourner une partie des nombres du jeton selon la
hauteur, une autre selon la largeur, et une troisième selon le temps. Le
principe ne change pas d'un pouce : il est simplement appliqué trois fois,
une fois par direction.

Les auteurs appellent cela un \emph{RoPE en trois dimensions}. Attention à
ne pas confondre ce mot avec la «~dimension du plongement~» vue plus haut :
ici, «~dimension~» veut dire «~direction dans laquelle on se repère~», et
rien d'autre.

\begin{enia}
Voici maintenant le point qui vous servira le jour où vous mettrez ce modèle
en service, et l'article y insiste. Comme la position se traduit par un
angle, on peut décider de \emph{resserrer} ou d'\emph{écarter} les crans.

Supposons que le modèle ait été entraîné sur des images de $16$ patchs de
côté, avec trente degrés par cran. On veut maintenant lui donner des images
de $32$ patchs de côté. Si l'on garde trente degrés par cran, la flèche fera
plusieurs tours complets et le modèle sera perdu. Mais si l'on passe à
quinze degrés par cran, la traversée complète de l'image correspondra
toujours au même angle total : le modèle retrouve ses repères.

Cette opération porte un nom, l'\emph{interpolation des fréquences}, et
l'article s'en sert vraiment, pour deux travaux que nous découvrirons au
chapitre~7 : la \emph{segmentation}, qui consiste à dire à quel objet
appartient chaque pixel, et l'\emph{estimation de profondeur}, qui consiste
à dire à quelle distance se trouve chaque pixel.

Ce que cela permet est énorme : \textbf{on peut donner au modèle des images
d'une autre taille que celles de son entraînement, sans le réentraîner}.
Nous y reviendrons au chapitre~26, quand nous choisirons nos réglages de
production.
\end{enia}

\begin{piege}
«~Sans le réentraîner~» ne veut pas dire «~sans rien vérifier~». Changer de
résolution modifie toujours un peu la qualité des résultats. La règle de
métier est simple : vous avez le droit de changer la résolution, mais vous
avez le devoir de mesurer ce que cela change sur vos propres données. Nous
verrons dans la partie~VII que l'article lui-même a mesuré cet effet, et
qu'il n'est pas nul.
\end{piege}

Nos jetons contiennent maintenant ce qu'ils voient et l'endroit d'où ils
viennent. Il est temps de les faire dialoguer.

\section{L'attention : chaque jeton écoute les autres}

\begin{enclair}
L'attention est le cœur du transformeur. Chaque jeton pose une question, tous
les autres présentent une étiquette, et celui qui pose la question écoute
surtout ceux dont l'étiquette répond le mieux. Il repart alors avec un
contenu enrichi de tout ce qu'il vient d'entendre.
\end{enclair}

Nous arrivons à la pièce maîtresse. Prenons le temps qu'il faut, car tout le
reste du livre s'appuie dessus.

Reprenons notre photo de chat. Le jeton qui contient les moustaches aimerait
bien savoir s'il y a des oreilles pointues quelque part dans l'image. Il ne
peut pas aller les chercher lui-même : il ne voit que son propre morceau. Il
va donc lancer un appel à toute la file, et voir qui répond.

Pour organiser cet appel, chaque jeton fabrique trois listes de nombres à
partir de son contenu, exactement comme il fabriquait son plongement : en
multipliant ses nombres par des coefficients appris. Ces trois listes ont
chacune un rôle bien précis.

\begin{definition}[La question, la clé et le contenu][qkv]
À partir de son propre jeton, chacun fabrique trois listes de nombres :

\begin{itemize}
  \item la \emph{question}, que l'on appelle \emph{query} en anglais : elle
        dit «~voici ce que je cherche~» ;
  \item la \emph{clé}, ou \emph{key} : elle dit «~voici de quoi je
        parle~» ;
  \item le \emph{contenu}, ou \emph{value} : c'est ce que le jeton donnera
        réellement à ceux qui l'auront écouté.
\end{itemize}

Les coefficients qui fabriquent ces trois listes sont appris pendant
l'entraînement. Personne ne les écrit à la main.
\end{definition}

\begin{intuition}
Imaginez une salle de classe. Un élève lève la main et demande : «~qui sait
quelque chose sur les oreilles pointues ?~» — c'est sa question. Chaque
autre élève porte une pancarte autour du cou où il a écrit son sujet : «~je
parle de fourrure~», «~je parle d'oreilles~», «~je parle de flaques d'eau~» —
ce sont les clés. L'élève qui a posé la question regarde toutes les
pancartes, repère celles qui collent le mieux à sa demande, et va écouter
ces camarades-là en priorité. Ce qu'ils lui racontent alors, c'est leur
contenu.
\end{intuition}

\begin{schema}{L'attention vue comme une salle de classe. Le jeton du bas
pose sa question. Chaque autre jeton présente sa clé. Plus la clé répond à
la question, plus la flèche est épaisse et plus le jeton sera écouté. Le
résultat est un mélange des contenus, dosé par l'épaisseur des flèches. Ici,
le jeton «~oreille~» est écouté fort, tandis que le jeton «~flaque~» est
presque ignoré.}[attention-classe]
  \begin{scope}[x=1cm, y=1cm]
    \node[bloc calcul, minimum width=20mm] (j1) at (0,3.2)   {fourrure};
    \node[bloc calcul, minimum width=20mm] (j2) at (3.2,3.2) {oreille};
    \node[bloc calcul, minimum width=20mm] (j3) at (6.4,3.2) {flaque};
    \node[font=\sffamily\scriptsize, text=mbSlate, above] at (0,3.72)
         {clé + contenu};
    \node[font=\sffamily\scriptsize, text=mbSlate, above] at (3.2,3.72)
         {clé + contenu};
    \node[font=\sffamily\scriptsize, text=mbSlate, above] at (6.4,3.72)
         {clé + contenu};

    \node[bloc accent, minimum width=30mm] (q) at (3.2,0) {moustaches};
    \node[font=\sffamily\scriptsize, text=mbAmber!80!black, below] at (3.2,-0.55)
         {pose la question : «~qui a des oreilles ?~»};

    \draw[draw=mbBlue!55, line width=1.1pt, -{Stealth[round,length=2.2mm]}]
          (j1.south) -- ($(q.north)+(-0.9,0)$);
    \draw[draw=mbBlue!85, line width=3.2pt, -{Stealth[round,length=3.2mm]}]
          (j2.south) -- (q.north);
    \draw[draw=mbBlue!30, line width=0.5pt, -{Stealth[round,length=1.8mm]}]
          (j3.south) -- ($(q.north)+(0.9,0)$);

    \node[etiquette] at (1.05,1.75) {$0{,}117$};
    \node[etiquette] at (3.62,1.75) {$0{,}867$};
    \node[etiquette] at (5.6,1.75) {$0{,}016$};
    \node[annotation] at (8.6,1.75) {les trois poids\\ font $1$ en tout};
  \end{scope}
\end{schema}

Comment mesure-t-on qu'une clé «~répond bien~» à une question ? On multiplie
les deux listes terme à terme et on additionne le tout. Cette opération porte
un nom en mathématiques : le \emph{produit scalaire}. Elle donne un grand
nombre quand les deux listes se ressemblent, et un petit nombre quand elles
n'ont rien à voir.

Le nombre obtenu s'appelle un \emph{score}. On en obtient un par jeton
écouté. Il faut ensuite transformer ces scores en \emph{poids},
c'est-à-dire en parts qui s'additionnent pour faire un, comme des parts d'un
gâteau. L'opération qui fait cela s'appelle la \emph{fonction softmax}.

\begin{definition}[La fonction softmax][softmax]
La \emph{softmax} transforme une liste de nombres quelconques en une liste de
parts, toutes positives, dont la somme fait exactement $1$. Elle procède en
deux temps : on calcule d'abord l'\emph{exponentielle} de chaque nombre, puis
on divise chacun par la somme de tous.

L'exponentielle est une machine à agrandir. On l'écrit $e^{x}$, ce qui se lit
«~e puissance $x$~». Elle prend un nombre et en rend un autre, toujours
positif, et qui grossit très vite : $e^{0} = 1$, $e^{1} \approx 2{,}7$,
$e^{2} \approx 7{,}4$, $e^{4} \approx 54{,}6$. Vous n'avez pas besoin de
savoir la calculer — toute calculatrice le fait. Retenez ses deux effets, qui
sont exactement ceux dont la softmax a besoin : elle rend tous les nombres
positifs, et elle écarte fortement les grands des petits.

Le nom «~softmax~» veut dire «~maximum en douceur~» : elle donne la plus
grosse part au plus grand score, sans pour autant écraser complètement les
autres.
\end{definition}

\begin{exemple}[Une attention calculée du début à la fin]
Faisons le calcul en entier, avec de tout petits nombres. Reprenons les
trois jetons de la figure~\ref{fig:attention-classe} : «~fourrure~»,
«~oreille~» et «~flaque~». Toutes les listes auront quatre nombres, et nous
conviendrons que le premier de ces quatre nombres veut dire «~je parle
d'oreilles~».

\textbf{Premier temps : la question.} Le jeton «~moustaches~» cherche des
oreilles. Il met donc du poids sur le premier nombre :
\[ q = (2 \ ; \ 0 \ ; \ 0 \ ; \ 0). \]

\textbf{Deuxième temps : les clés.} Les trois autres présentent les leurs.
«~Oreille~» parle beaucoup d'oreilles, «~fourrure~» un peu, «~flaque~» pas
du tout :
\[
  \begin{aligned}
    k_{\text{fourrure}} &= (2 \ ; \ 0 \ ; \ 0 \ ; \ 0), \\
    k_{\text{oreille}}  &= (4 \ ; \ 0 \ ; \ 0 \ ; \ 0), \\
    k_{\text{flaque}}   &= (0 \ ; \ 3 \ ; \ 0 \ ; \ 0).
  \end{aligned}
\]

\textbf{Troisième temps : les scores.} On fait le produit scalaire de la
question avec chaque clé :
\[
  2 \times 2 = 4, \qquad
  2 \times 4 = 8, \qquad
  0 .
\]

\textbf{Quatrième temps : on remet à l'échelle.} On divise chaque score par
la racine carrée de la longueur des listes, ici $\sqrt{4} = 2$. La racine
carrée
d'un nombre, c'est celui qui, multiplié par lui-même, redonne ce nombre :
$\sqrt{4} = 2$ parce que $2 \times 2 = 4$, et $\sqrt{9} = 3$ parce que
$3 \times 3 = 9$. Cette division empêche les scores de devenir énormes quand
les listes sont longues :
\[
  \frac{4}{2} = 2, \qquad \frac{8}{2} = 4, \qquad \frac{0}{2} = 0 .
\]

\textbf{Cinquième temps : la softmax.} On prend l'exponentielle de chacun :
\[
  e^{2} \approx 7{,}389, \qquad e^{4} \approx 54{,}598, \qquad e^{0} = 1 .
\]
On les additionne : $7{,}389 + 54{,}598 + 1 = 62{,}987$. Puis on divise
chacun par cette somme :
\[
  \frac{7{,}389}{62{,}987} \approx 0{,}117, \quad
  \frac{54{,}598}{62{,}987} \approx 0{,}867, \quad
  \frac{1}{62{,}987} \approx 0{,}016 .
\]
Vérifions : $0{,}117 + 0{,}867 + 0{,}016 = 1{,}000$. Ce sont bien des parts
de gâteau, et ce sont exactement les trois nombres écrits sur les flèches de
la figure~\ref{fig:attention-classe}.

\textbf{Sixième temps : le mélange des contenus.} Disons que les trois
contenus valent
\[
  \begin{aligned}
    v_{\text{fourrure}} &= (10 \ ; \ 0 \ ; \ 0 \ ; \ 0), \\
    v_{\text{oreille}}  &= (0 \ ; \ 10 \ ; \ 0 \ ; \ 0), \\
    v_{\text{flaque}}   &= (0 \ ; \ 0 \ ; \ 10 \ ; \ 0).
  \end{aligned}
\]
Le résultat est la somme des contenus, chacun multiplié par sa part :
\[
  \begin{aligned}
    &0{,}117 \times v_{\text{fourrure}}
     + 0{,}867 \times v_{\text{oreille}}
     + 0{,}016 \times v_{\text{flaque}} \\
    &\qquad = (1{,}17 \ ; \ 8{,}67 \ ; \ 0{,}16 \ ; \ 0).
  \end{aligned}
\]

Voilà le nouveau jeton «~moustaches~». Il ne contient plus seulement ce
qu'il voyait tout seul : il contient surtout ce que lui a dit «~oreille~»,
un peu de ce que lui a dit «~fourrure~», et presque rien de «~flaque~». Il a
posé une question sur les oreilles, et c'est bien le jeton «~oreille~» qu'il
a le plus écouté.
\end{exemple}

\begin{piege}
N'oubliez jamais la division par la racine carrée. Sans elle, avec des
listes de mille nombres, les scores atteignent des valeurs si grandes que la
softmax donne une part de $1$ à un seul jeton et $0$ à tous les autres. Le
modèle n'écoute plus qu'un seul camarade, l'entraînement se bloque, et rien
dans les messages d'erreur ne vous dira pourquoi.
\end{piege}

\begin{aretenir}
L'attention se résume à six gestes, toujours les mêmes : je fabrique ma
question, chacun fabrique sa clé et son contenu, je compare ma question à
chaque clé, je remets à l'échelle, je transforme les scores en parts, je
mélange les contenus selon ces parts. Tout ce qui suit dans ce chapitre
n'est que de l'organisation autour de ces six gestes.
\end{aretenir}

Cette façon de faire est puissante, mais il lui manque quelque chose. Voyons
quoi, et comment on y remédie.

\section{Plusieurs têtes valent mieux qu'une}

\begin{enclair}
Avec une seule attention, chaque jeton ne peut poser qu'une seule question à
la fois. On en fait donc tourner plusieurs en même temps, chacune avec
ses propres réglages, et on met leurs réponses bout à bout. Chacune de ces
attentions qui travaillent côte à côte s'appelle une tête.
\end{enclair}

Notre jeton «~moustaches~» cherchait des oreilles. Mais il aurait aussi aimé
savoir quelle est la couleur dominante autour de lui, et si quelque chose
bouge à sa gauche. Avec une seule question, il doit choisir. C'est dommage.

\begin{definition}[La tête d'attention][tete-attention]
Une \emph{tête d'attention} est un exemplaire complet de tout ce que nous
venons de décrire, avec ses propres coefficients pour fabriquer les
questions, les clés et les contenus. Plusieurs têtes travaillent en même
temps sur les mêmes jetons, mais chacune pose sa propre question. On parle
alors d'\emph{attention multi-têtes}.
\end{definition}

Le partage se fait ainsi : on ne donne pas la totalité des nombres du jeton
à chaque tête, on les répartit entre elles. Chaque tête travaille sur sa
part, produit sa réponse, et l'on recolle toutes les réponses bout à bout à
la fin. La liste finale a donc la même longueur qu'au départ.

\begin{schema}{Le partage entre les têtes. Les nombres du jeton sont répartis
en parts égales, une par tête. Chaque tête pose sa propre question et ne
travaille que sur sa part. On recolle ensuite les réponses bout à bout, et
l'on retrouve une liste de la même longueur qu'au départ. Quatre têtes sont
dessinées ici ; le modèle en emploie seize.}[tetes-attention]
  \begin{scope}[x=1cm, y=1cm]
    \node[font=\sffamily\scriptsize, text=mbSlate] at (-1.5,3.3) {le jeton};
    \foreach \k in {0,...,3}{
      \pgfmathtruncatemacro{\tone}{22+16*\k}
      \fill[mbAzure!\tone] (\k*1.9,3.05) rectangle ++(1.9,0.5);
      \draw[draw=white, line width=1.6pt] (\k*1.9,3.05) rectangle ++(1.9,0.5);
      \draw[draw=mbSlate!50, line width=0.5pt] (\k*1.9,3.05) rectangle ++(1.9,0.5);
    }

    \foreach \k/\quoi in {0/{les formes}, 1/{les couleurs},
                          2/{les mouvements}, 3/{les bords}}{
      \draw[fleche] (\k*1.9+0.95,3.0) -- (\k*1.9+0.95,2.05);
      \node[bloc calcul, minimum width=17mm, minimum height=8mm,
            font=\sffamily\scriptsize] (t\k) at (\k*1.9+0.95,1.6) {tête \the\numexpr\k+1\relax};
      \node[font=\sffamily\scriptsize, text=mbAmber!80!black]
            at (\k*1.9+0.95,0.85) {\quoi};
      \draw[fleche] (\k*1.9+0.95,0.55) -- (\k*1.9+0.95,-0.2);
    }

    \foreach \k in {0,...,3}{
      \pgfmathtruncatemacro{\tone}{22+16*\k}
      \fill[mbTeal!\tone] (\k*1.9,-0.75) rectangle ++(1.9,0.5);
      \draw[draw=white, line width=1.6pt] (\k*1.9,-0.75) rectangle ++(1.9,0.5);
      \draw[draw=mbSlate!50, line width=0.5pt] (\k*1.9,-0.75) rectangle ++(1.9,0.5);
    }
    \node[font=\sffamily\scriptsize, text=mbSlate] at (-1.6,-0.5)
         {on recolle};
    \draw[accolade bas] (0,-0.9) -- (7.6,-0.9);
    \node[font=\sffamily\scriptsize, text=mbSlate] at (3.8,-1.55)
         {la même longueur qu'au départ};
  \end{scope}
\end{schema}

\begin{exemple}[Répartir mille vingt-quatre nombres entre seize têtes]
Prenons un jeton de $1\,024$ nombres et $16$ têtes d'attention. Chaque tête
reçoit
\[
  \frac{1\,024}{16} = 64 \text{ nombres.}
\]
Chacune produit une réponse de $64$ nombres. En recollant les seize
réponses, on retrouve
\[
  16 \times 64 = 1\,024 \text{ nombres.}
\]
La taille du jeton n'a pas changé, mais seize questions différentes ont été
posées au lieu d'une seule, et pour à peu près le même prix.
\end{exemple}

\begin{intuition}
Retournons dans la salle de classe. Avec une seule tête, l'élève ne peut
poser qu'une question et doit attendre la réponse avant d'en poser une
autre. Avec seize têtes, c'est comme s'il avait seize paires d'oreilles :
l'une écoute les formes, l'autre les couleurs, une troisième les
mouvements. Il rassemble ensuite tout ce qu'il a entendu en une seule fiche.
\end{intuition}

Nous tenons la pièce centrale. Il faut maintenant l'emballer proprement, et
c'est le rôle du bloc de transformeur.

\section{Le bloc de transformeur, ouvert en coupe}

\begin{enclair}
Un bloc de transformeur, c'est l'attention plus trois pièces qui la rendent
utilisable : une remise au propre des nombres avant chaque étage, un petit
réseau qui réfléchit sur chaque jeton séparément, et deux raccourcis qui
laissent passer l'information sans la déformer.
\end{enclair}

L'attention toute seule ne suffit pas. Elle mélange les jetons entre eux,
mais elle ne fait aucun raisonnement à l'intérieur d'un jeton. Et surtout,
si l'on empile beaucoup d'attentions à la suite, les nombres partent dans
tous les sens et l'entraînement échoue. Trois pièces règlent ces problèmes.

\subsection*{Première pièce : la normalisation de couche}

\begin{definition}[La normalisation de couche, ou LayerNorm][layernorm]
La \emph{normalisation de couche} remet les nombres d'un jeton à une échelle
raisonnable. Pour chaque jeton pris séparément, on calcule la moyenne de ses
nombres, on la retire de chacun, puis on divise le tout par un nombre qui
mesure à quel point ils s'écartaient de cette moyenne. Après cette
opération, les nombres du jeton sont centrés autour de zéro et ne sont ni
minuscules ni gigantesques. On l'appelle presque toujours par son nom
anglais, \emph{LayerNorm}.
\end{definition}

\begin{exemple}[Normaliser quatre nombres]
Un jeton contient les nombres $(2 \ ; \ 4 \ ; \ 6 \ ; \ 8)$.

Sa moyenne vaut $(2+4+6+8)/4 = 20/4 = 5$. On la retire de chacun :
\[ (-3 \ ; \ -1 \ ; \ 1 \ ; \ 3). \]
On mesure ensuite de combien les nombres s'écartent en moyenne. On multiplie
chaque écart par lui-même — c'est ce qu'on appelle l'élever au carré — ce
qui donne $9$, $1$, $1$, $9$. Multiplier par soi-même sert ici à effacer les
signes moins, car un nombre négatif multiplié par lui-même devient positif.
La moyenne de ces carrés vaut $(9+1+1+9)/4 = 5$, et sa racine carrée vaut
$\sqrt{5} \approx 2{,}236$. On divise :
\[
  \left( \frac{-3}{2{,}236} \ ; \ \frac{-1}{2{,}236} \ ; \
         \frac{1}{2{,}236} \ ; \ \frac{3}{2{,}236} \right)
  \approx (-1{,}342 \ ; \ -0{,}447 \ ; \ 0{,}447 \ ; \ 1{,}342).
\]

Les nombres tiennent maintenant à peu près entre $-1{,}5$ et $1{,}5$. Le
jeton dit toujours la même chose — les rapports entre ses nombres n'ont pas
changé — mais il le dit à voix normale au lieu de crier ou de chuchoter.
\end{exemple}

Voilà nos nombres remis d'aplomb. Passons à la pièce qui va s'en servir.

\subsection*{Deuxième pièce : le petit réseau de réflexion}

Après avoir écouté les autres, chaque jeton a besoin d'un moment pour
digérer ce qu'il vient d'entendre. C'est le rôle du second étage.

\begin{definition}[Le MLP][mlp]
Un \emph{MLP} est un petit réseau appliqué à chaque jeton séparément, sans
aucun regard sur les autres. Ces trois lettres viennent de l'anglais
\emph{Multi-Layer Perceptron}. Ce nom ne vous servira jamais : retenez
seulement «~petit réseau de réflexion~». Il procède en
trois temps : il agrandit d'abord la liste de nombres, en général en la
multipliant par quatre ; il applique ensuite à chaque nombre une petite
opération de tri appelée \emph{fonction d'activation} ; enfin il ramène la
liste à sa taille de départ.
\end{definition}

\begin{definition}[La fonction d'activation][activation]
Vous connaissez déjà cette pièce, sans connaître son nom. Dans le programme
du chapitre~\ref{chap:auto-supervision}, nous avions glissé une petite règle
appelée \co{python}{ReLU}, qui remplaçait par zéro tout nombre négatif et
laissait passer les autres. Nous avions dit alors que, sans elle, empiler
deux couches ne servirait à rien.

Cette règle porte un nom de famille : on l'appelle une \emph{fonction
d'activation}. C'est une opération appliquée à chaque nombre séparément, et
qui ne se contente pas de le multiplier ou de lui ajouter quelque chose. La
\co{python}{ReLU} est la plus simple de la famille ; celle qu'emploient les
transformeurs, la \co{python}{GELU}, arrondit un peu les angles mais rend le
même service. Sans l'une d'elles, toute la pile de blocs se ramènerait à une
seule multiplication suivie d'une seule addition.
\end{definition}

\begin{exemple}[Un MLP en tout petit]
Prenons un jeton de deux nombres, $(3 \ ; \ -2)$, et un MLP qui agrandit la
liste en la multipliant par deux, donc qui passe de deux à quatre nombres.

\textbf{On agrandit.} Après le premier calcul, disons que l'on obtient
\[ (4 \ ; \ -1 \ ; \ 6 \ ; \ -5). \]

\textbf{On trie.} La fonction d'activation remplace les négatifs par zéro :
\[ (4 \ ; \ 0 \ ; \ 6 \ ; \ 0). \]

Regardez ce qui vient de se passer : deux des quatre mesures ont été
éteintes, deux sont restées allumées. Le jeton ne réagit donc qu'à certaines
situations et reste muet devant les autres. C'est exactement ce qu'une suite
de multiplications ne saurait pas faire.

\textbf{On revient à la taille de départ.} Un dernier calcul ramène ces
quatre nombres à deux, et le résultat est ajouté au jeton d'origine — nous
allons voir tout de suite pourquoi.
\end{exemple}

\begin{intuition}
L'attention, c'est le moment où l'on discute avec ses voisins. Le MLP, c'est
le moment où l'on retourne à sa place pour réfléchir tout seul à ce qu'on
vient d'apprendre. Un bloc de transformeur fait donc alterner, encore et
encore, un temps de discussion et un temps de réflexion personnelle.
\end{intuition}

Il nous reste une pièce, la plus discrète des trois, et pourtant celle sans
laquelle rien ne tiendrait debout.

\subsection*{Troisième pièce : la connexion résiduelle}

\begin{definition}[La connexion résiduelle][residuelle]
Une \emph{connexion résiduelle} est un raccourci : au lieu de remplacer le
jeton par le résultat d'un étage, on \emph{ajoute} ce résultat au jeton
d'origine. L'ancienne valeur n'est donc jamais perdue ; l'étage ne fait
qu'apporter une correction. On dit aussi \emph{connexion saute-mouton},
parce que l'information peut sauter par-dessus l'étage.
\end{definition}

\begin{exemple}[Ce que le raccourci change, en chiffres]
Un jeton vaut $(3 \ ; \ -2)$ et l'étage qu'il traverse produit
$(0{,}4 \ ; \ 0{,}1)$.

\textbf{Sans raccourci}, le jeton devient $(0{,}4 \ ; \ 0{,}1)$ : tout ce
qu'il contenait a disparu, remplacé par le travail de l'étage.

\textbf{Avec raccourci}, on additionne :
\[ (3 \ ; \ -2) + (0{,}4 \ ; \ 0{,}1) = (3{,}4 \ ; \ -1{,}9). \]

Le jeton est toujours reconnaissable, simplement corrigé un peu. Imaginez
maintenant quarante-huit étages à la suite : dans le premier cas, il ne reste
plus rien du jeton de départ après quelques étages ; dans le second, il est
encore là, augmenté de quarante-huit petites corrections.
\end{exemple}

\begin{intuition}
Vous relisez le devoir d'un élève. Vous pouvez soit tout réécrire de votre
main — et si vous vous trompez, le travail de l'élève est perdu — soit
écrire vos remarques en marge, à côté de son texte. La connexion résiduelle,
c'est la deuxième façon de faire. Elle explique pourquoi l'on peut empiler
quarante-huit blocs les uns sur les autres sans que le message du départ se
dissolve en chemin.
\end{intuition}

Assemblons maintenant les trois pièces autour de l'attention.

\begin{schema}{Un bloc de transformeur ouvert en coupe. Le jeton entre par
le bas. Il passe dans deux étages successifs. Chaque étage commence par une
normalisation, puis fait son travail — écouter les autres pour le premier,
réfléchir tout seul pour le second — et son résultat est \emph{ajouté} à ce
qui entrait dans l'étage, grâce au raccourci dessiné sur le côté. Le jeton
ressort par le haut, de la même taille qu'à l'entrée.}[bloc-transformeur]
  \begin{scope}[x=1cm, y=1cm]
    \node[font=\sffamily\small\bfseries, text=mbSlate] at (0,-0.85)
         {jeton qui entre};
    \draw[fleche] (0,-0.55) -- (0,0.15);

    %---- premier étage : le jeton écoute les autres ------------------
    \node[bloc, minimum width=48mm, text width=42mm] (n1) at (0,0.7)
         {normalisation};
    \node[bloc calcul, minimum width=48mm, text width=42mm] (a1) at (0,2.3)
         {attention multi-têtes};
    \node[circle, draw=mbAmber!85!black, fill=mbAmber!12, line width=0.9pt,
          minimum size=8mm, inner sep=0pt, font=\small] (s1) at (0,3.9) {$+$};

    \draw[fleche] (n1) -- (a1);
    \draw[fleche] (a1) -- (s1);
    \draw[fleche accent] (-3.5,0.15) -- (-3.5,3.9) -- (s1);
    \draw[draw=mbSlate!45, line width=0.6pt] (0,0.15) -- (-3.5,0.15);
    \node[annotation, rotate=90] at (-3.82,2.05) {raccourci};

    %---- second étage : le jeton réfléchit tout seul -----------------
    \node[bloc, minimum width=48mm, text width=42mm] (n2) at (0,5.3)
         {normalisation};
    \node[bloc calcul, minimum width=48mm, text width=42mm] (a2) at (0,6.9)
         {petit réseau de réflexion};
    \node[circle, draw=mbAmber!85!black, fill=mbAmber!12, line width=0.9pt,
          minimum size=8mm, inner sep=0pt, font=\small] (s2) at (0,8.5) {$+$};

    \draw[fleche] (s1) -- (n2);
    \draw[fleche] (n2) -- (a2);
    \draw[fleche] (a2) -- (s2);
    \draw[fleche accent] (-3.5,4.65) -- (-3.5,8.5) -- (s2);
    \draw[draw=mbSlate!45, line width=0.6pt] (0,4.65) -- (-3.5,4.65);
    \node[annotation, rotate=90] at (-3.82,6.6) {raccourci};

    \draw[fleche] (s2) -- (0,9.3);
    \node[font=\sffamily\small\bfseries, text=mbSlate] at (0,9.6)
         {jeton qui sort};

    %---- les bulles, posées à distance des cadres --------------------
    \node[annotation bulle, text width=2.9cm, anchor=west] at (3.1,2.3)
         {ici, le jeton\\ discute avec\\ les autres};
    \node[annotation bulle, text width=2.9cm, anchor=west] at (3.1,6.9)
         {ici, il réfléchit\\ tout seul};
  \end{scope}
\end{schema}

\begin{definition}[Le bloc de transformeur][bloc-transformeur]
Un \emph{bloc de transformeur} est l'ensemble que nous venons d'assembler :
une normalisation puis une attention multi-têtes, le tout ajouté à l'entrée ;
puis une normalisation puis un petit réseau de réflexion, le tout ajouté de
nouveau. Un bloc reçoit une file de jetons et rend une file de jetons de la
même taille. C'est pour cela qu'on peut en empiler autant qu'on veut.
\end{definition}

\begin{exemple}[Un bloc, en chiffres]
Donnons à un bloc les $2\,048$ jetons d'un clip de la phase principale,
chacun fait de $1\,024$ nombres. À la sortie, nous récupérons encore
$2\,048$ jetons de $1\,024$ nombres.

Rien n'a changé en taille, tout a changé en contenu : chaque jeton a écouté
les $2\,047$ autres, puis a réfléchi de son côté. Comme l'entrée et la
sortie ont exactement la même forme, on peut brancher la sortie de ce bloc
sur l'entrée du suivant, sans le moindre ajustement. C'est ce détail, et lui
seul, qui rend les empilements possibles.
\end{exemple}

Et c'est justement ce que l'on fait dans la section suivante.

\section{Empiler les blocs : la profondeur et les quatre tailles}

\begin{enclair}
On met les blocs les uns au-dessus des autres, et la file de jetons les
traverse tous. Le nombre de blocs empilés s'appelle la profondeur. Plus on
monte, plus ce que voient les jetons devient général.
\end{enclair}

\begin{definition}[La profondeur][profondeur]
La \emph{profondeur} d'un transformeur est le nombre de blocs empilés les uns
sur les autres. Les jetons traversent le premier bloc, puis le deuxième, et
ainsi de suite jusqu'au dernier. Chaque bloc reçoit le travail du précédent.
\end{definition}

Ce qui se passe le long de cette pile vaut la peine qu'on s'y arrête, et
nous nous en servirons énormément dans la partie~IV. Les premiers blocs
repèrent des choses très simples : un bord, un changement de couleur, un
léger déplacement. Les blocs du milieu assemblent ces petites choses en
morceaux reconnaissables : un coin, une roue, une main. Les derniers blocs
assemblent ces morceaux en idées complètes : une voiture, quelqu'un qui
verse de l'eau.

\begin{schema}{Ce que voient les blocs, du bas vers le haut. En bas, on
repère des traits. Un peu plus haut, des coins. Plus haut encore, des
morceaux d'objets. Tout en haut, l'objet entier et l'action. C'est cette
montée en généralité qui explique pourquoi, plus tard, nous aurons parfois
envie d'aller chercher des informations à mi-hauteur plutôt qu'au
sommet.}[montee-blocs]
  \begin{scope}[x=1cm, y=1cm]
    \foreach \k/\nom/\quoi in {
      0/{blocs du bas}/{des traits},
      1/{blocs du milieu}/{des coins},
      2/{blocs du haut}/{des roues, des mains},
      3/{dernier bloc}/{une voiture, un geste}}{
      \pgfmathsetmacro{\yy}{\k*1.35}
      \pgfmathtruncatemacro{\tone}{12+16*\k}
      \node[rectangle, rounded corners=2.5pt, draw=mbBlue, fill=mbBlue!\tone,
            line width=0.8pt, minimum width=38mm, minimum height=9mm,
            font=\sffamily\small, text=mbInk] (b\k) at (0,\yy) {\nom};
      \node[font=\sffamily\small, text=mbSlate, anchor=west] at (2.4,\yy)
           {il y voit \quoi};
    }
    \draw[fleche] (b0) -- (b1);
    \draw[fleche] (b1) -- (b2);
    \draw[fleche] (b2) -- (b3);
    \draw[fleche accent] (-2.5,-0.55) -- (-2.5,4.6);
    \node[annotation, rotate=90] at (-2.85,2.0) {de plus en plus général};
  \end{scope}
\end{schema}

\begin{piege}
Cette montée en généralité a un revers, et c'est exactement le problème que
V-JEPA~2.1 vient résoudre. Si le sommet ne parle plus que d'idées générales,
il a pu oublier \emph{où} les choses se trouvaient. Or pour dire à quelle
distance est chaque pixel, on a justement besoin du «~où~». Gardez cette
tension en tête : elle est le sujet de toute la partie~III.
\end{piege}

\begin{piege}
Il faut maintenant corriger une image que nous vous avions donnée au
chapitre~\ref{chap:auto-supervision}. Nous y avions dessiné l'encodeur comme
un \emph{entonnoir} : beaucoup de nombres entrent, un petit paquet sort. Cette
image était commode pour comprendre l'idée de résumé, mais elle est trop
simple.

La vérité, vous venez de la voir : l'encodeur ne rend pas \emph{un} paquet
pour toute l'image. Il rend \emph{un paquet par jeton}. Un clip de la phase
principale entre avec $2\,048$ jetons, et il ressort avec $2\,048$ jetons,
un pour chaque endroit du clip. L'entonnoir n'a pas rétréci la file : il a
rétréci chaque morceau, chacun de son côté.

Retenez bien cette différence, car tout le livre repose dessus. S'il vous
faut une seule réponse pour toute la vidéo — « que se passe-t-il ici ? » —
il faudra mélanger ces $2\,048$ paquets en un seul. Mais s'il vous faut une
réponse \emph{par endroit} — « à quelle distance est ce pixel-là ? » — il
faudra que chaque paquet, pris tout seul, décrive bien son propre morceau.
Ces deux besoins tirent le modèle dans des directions différentes, et c'est
précisément le conflit que raconte la partie~III.
\end{piege}

Combien de blocs empile-t-on, et de quelle taille ? La communauté s'est mise
d'accord sur quelques tailles standards, désignées par une lettre. V-JEPA~2.1
en emploie quatre.

\begin{surlesdonnees}
Les quatre tailles employées dans l'article, de la plus petite à la plus
grande :

\begin{center}
\footnotesize
\begin{tabular}{llrl}
\toprule
\textbf{Nom} & \textbf{Se lit} & \textbf{Boutons} & \textbf{Rôle dans l'article} \\
\midrule
ViT-B & «~bé~», pour \emph{Base}     & $80$ M            & le petit modèle de service \\
ViT-L & «~èl~», pour \emph{Large}    & $300$ M           & le modèle des essais \\
ViT-g & «~gé~» minuscule, \emph{giant} & $1$ milliard    & la taille de V-JEPA~2 \\
ViT-G & «~gé~» majuscule, \emph{Giant} & $2$ milliards   & le grand modèle final \\
\bottomrule
\end{tabular}
\end{center}

Les deux plus petits ne sont pas entraînés depuis le début : ils apprennent
en copiant le grand, par un procédé appelé \emph{distillation} que nous
verrons au chapitre~15. Retenez seulement, pour l'instant, qu'un petit modèle
peut hériter de presque toute la qualité d'un grand.

Attention à la lettre \textbf{g}~: la minuscule et la majuscule désignent
deux modèles différents, et l'article emploie les deux à quelques lignes
d'écart. Le grand modèle final, celui qui obtient tous les meilleurs
résultats, est le ViT-G majuscule et compte $48$ blocs empilés. Nous le
savons parce que l'article indique qu'il va chercher des informations aux
blocs numéro $12$, $24$, $36$ et $48$ — nous verrons au
chapitre~11 pourquoi il fait cela.
\end{surlesdonnees}

\begin{exemple}[Ce que «~deux milliards de boutons~» veut dire]
Reprenez le chapitre~\ref{chap:apprendre} : un bouton, c'est un nombre que
l'entraînement règle. Le ViT-G en possède deux milliards. Si vous les
imprimiez à raison de cinquante nombres par ligne et cinquante lignes par
page, il vous faudrait huit cent mille pages, soit une pile de papier de
quatre-vingts mètres de haut. Et l'entraînement les règle \emph{tous},
$147\,000$ fois de suite.

C'est aussi pour cela que nous consacrerons tout le chapitre~15 à rétrécir
ce modèle : personne ne veut faire tourner quatre-vingts mètres de papier
dans un serveur, et l'article montre comment ramener presque toute cette
qualité dans un modèle six fois plus petit.
\end{exemple}

Il nous reste une dernière chose à regarder en face avant de coder : le prix
de tout cela.

\section{Ce que l'attention coûte vraiment}

\begin{enclair}
Dans un transformeur, le nombre de jetons est le bouton le plus important
pour la vitesse. Le doubler ne double pas le travail : il le multiplie par
quatre. C'est la première chose à regarder quand on choisit une résolution.
\end{enclair}

Nous l'avons dit au début du chapitre, mais nous pouvons maintenant le dire
avec de vrais chiffres. Dans l'attention, chaque jeton compare sa question à
la clé de tous les autres. Avec $n$ jetons, cela fait $n \times n$
comparaisons.

\begin{exemple}[Le prix des deux configurations de V-JEPA 2.1]
Phase principale, $2\,048$ jetons :
\[ 2\,048 \times 2\,048 = 4\,194\,304 \text{ comparaisons.} \]

Phase de refroidissement, $18\,432$ jetons :
\[ 18\,432 \times 18\,432 = 339\,738\,624 \text{ comparaisons.} \]

Le rapport vaut
\[ \frac{339\,738\,624}{4\,194\,304} = 81 . \]

Quatre-vingt-une fois plus de travail, et cela dans \emph{chacun} des
quarante-huit blocs. Voilà le chiffre qui explique tout le calendrier de
l'entraînement décrit dans l'article.
\end{exemple}

\begin{aretenir}
Retenez la règle du praticien : \textbf{ce que vous paierez suit à peu près
le carré du nombre de jetons}. Quand votre modèle sera trop lent
en service, le premier bouton à tourner ne sera ni le processeur ni la
bibliothèque : ce sera la résolution d'entrée, parce que c'est elle qui
décide du nombre de jetons. Nous referons ce calcul, pour de vrai, au
chapitre~29.
\end{aretenir}

Nous avons monté l'œil de la machine en entier. Vérifions maintenant, avec
un programme, que tout ce que nous avons calculé à la main est bien ce qui se
passe pour de vrai.

\section{Passons au code}

Nous allons écrire quatre petits programmes. Les deux premiers découpent une
image et un clip en jetons, et retrouvent les quatre nombres du tableau. Le
troisième refait, ligne par ligne, l'exemple d'attention calculé à la main.
Le quatrième assemble un bloc de transformeur complet. Vous pouvez les taper
tels quels : ils tournent sur n'importe quel ordinateur, même sans le
processeur spécialisé dont les gros entraînements ont besoin.

Commençons par le découpage. Comme annoncé, on l'obtient avec une
convolution dont la fenêtre et le pas ont la même taille.

\begin{code}{python}[Découper une image et un clip en jetons]
import torch
import torch.nn as nn

# --- La voie image : une convolution 2D de 16 x 16, qui avance de 16 ---
tokeniseur_image = nn.Conv2d(
    in_channels=3,       # les trois couleurs
    out_channels=1024,   # la taille d'un jeton
    kernel_size=16,      # la fenetre : un patch de 16 x 16
    stride=16,           # elle avance de 16 : aucun chevauchement
)

image = torch.zeros(1, 3, 256, 256)          # (lot, couleurs, hauteur, largeur)
sortie = tokeniseur_image(image)             # -> (1, 1024, 16, 16)
jetons_image = sortie.flatten(2).transpose(1, 2)

print(jetons_image.shape)   # -> torch.Size([1, 256, 1024])
# 256 jetons : c'est bien le nombre calcule a la main.

# --- La voie video : une convolution 3D, avec 2 images d'epaisseur ---
tokeniseur_video = nn.Conv3d(
    in_channels=3,
    out_channels=1024,
    kernel_size=(2, 16, 16),   # (tubelet, hauteur, largeur)
    stride=(2, 16, 16),
)

clip = torch.zeros(1, 3, 16, 256, 256)       # (lot, couleurs, temps, H, W)
sortie = tokeniseur_video(clip)              # -> (1, 1024, 8, 16, 16)
jetons_video = sortie.flatten(2).transpose(1, 2)

print(jetons_video.shape)   # -> torch.Size([1, 2048, 1024])
# 2048 jetons : 8 x 16 x 16, exactement notre calcul.
\end{code}

\begin{piege}
Regardez bien l'ordre des axes : PyTorch attend les couleurs en deuxième
position, pas en dernière. C'est le changement d'ordre que nous avions déjà
rencontré au chapitre~\ref{chap:image-video}. Si vous lui donnez un tenseur
rangé dans l'ordre de l'affichage, il ne se plaindra pas toujours : il
calculera quelque chose de faux, en silence. Prenez l'habitude d'afficher la
forme avant et après chaque étape.
\end{piege}

Vérifions maintenant les quatre nombres du tableau, tous d'un coup.

\begin{code}{python}[Vérifier les quatre nombres de jetons]
def compte_jetons(images, hauteur, largeur, patch=16, tubelet=2):
    """Rend le nombre de jetons. Mettre images=None pour une image seule."""
    n = (hauteur // patch) * (largeur // patch)
    if images is not None:
        n = n * (images // tubelet)
    return n

print(compte_jetons(None, 256, 256))   # -> 256
print(compte_jetons(None, 512, 512))   # -> 1024
print(compte_jetons(16,  256, 256))    # -> 2048
print(compte_jetons(64,  384, 384))    # -> 18432

# Et le prix de l'attention, qui grimpe au carre :
print(18432**2 / 2048**2)              # -> 81.0
\end{code}

Passons à l'attention. Nous reprenons exactement les nombres de l'exemple
calculé à la main : vous devez retrouver les mêmes parts, à la troisième
décimale.

\begin{code}{python}[L'attention, refaite avec nos nombres]
import torch

# Notre question, puis les cles et les contenus de fourrure, oreille, flaque.
q = torch.tensor([2.0, 0.0, 0.0, 0.0])
k = torch.tensor([[2.0, 0.0, 0.0, 0.0],    # fourrure
                  [4.0, 0.0, 0.0, 0.0],    # oreille
                  [0.0, 3.0, 0.0, 0.0]])   # flaque
v = torch.tensor([[10.0,  0.0,  0.0, 0.0],
                  [ 0.0, 10.0,  0.0, 0.0],
                  [ 0.0,  0.0, 10.0, 0.0]])

scores = k @ q                      # les produits scalaires
print(scores)                       # -> tensor([4., 8., 0.])

scores = scores / (q.shape[0] ** 0.5)   # on divise par racine de 4, donc par 2
print(scores)                       # -> tensor([2., 4., 0.])

poids = torch.softmax(scores, dim=0)
print(poids)                        # -> tensor([0.1173, 0.8668, 0.0159])
print(poids.sum())                  # -> tensor(1.)

resultat = poids @ v
print(resultat)                     # -> tensor([1.1731, 8.6681, 0.1588, 0.])
\end{code}
Les trois parts sont bien $0{,}117$, $0{,}867$ et $0{,}016$, dans cet ordre,
et le résultat est bien celui que nous avions trouvé à la main. Le programme
ne fait rien d'autre que ce que nous avons écrit sur le papier.

Terminons par un bloc de transformeur complet. Comptez les lignes : il n'y
en a pas vingt. Tout ce chapitre tient là-dedans.

\begin{code}{python}[Un bloc de transformeur, en entier]
import torch
import torch.nn as nn

class BlocTransformeur(nn.Module):
    def __init__(self, dimension=1024, nb_tetes=16):
        super().__init__()
        self.norme1 = nn.LayerNorm(dimension)
        self.attention = nn.MultiheadAttention(
            dimension, nb_tetes, batch_first=True)
        self.norme2 = nn.LayerNorm(dimension)
        self.reflexion = nn.Sequential(
            nn.Linear(dimension, 4 * dimension),   # on agrandit
            nn.GELU(),                             # la fonction d'activation
            nn.Linear(4 * dimension, dimension),   # on revient a la taille de depart
        )

    def forward(self, jetons):
        # Premier etage : on ecoute les autres, et on AJOUTE le resultat.
        propre = self.norme1(jetons)
        ecoute, _ = self.attention(propre, propre, propre)
        jetons = jetons + ecoute

        # Second etage : on reflechit tout seul, et on AJOUTE encore.
        jetons = jetons + self.reflexion(self.norme2(jetons))
        return jetons

bloc = BlocTransformeur()
jetons = torch.zeros(1, 2048, 1024)     # un clip de la phase principale
print(bloc(jetons).shape)               # -> torch.Size([1, 2048, 1024])

# La taille ne change pas : c'est ce qui permet d'empiler 48 blocs.
encodeur = nn.Sequential(*[BlocTransformeur() for _ in range(48)])
\end{code}

\begin{piege}
Repérez les deux \co{python}{jetons = jetons + ...} : ce sont les connexions
résiduelles. Si vous écrivez \co{python}{jetons = ecoute} au lieu de
\co{python}{jetons = jetons + ecoute}, le programme tournera sans erreur,
mais l'entraînement d'un modèle profond ne convergera jamais. C'est l'une
des erreurs les plus fréquentes, et l'une des plus difficiles à repérer,
parce que rien ne casse : les résultats sont simplement mauvais.
\end{piege}

\begin{aretenir}
Un transformeur de vision, c'est quatre gestes et rien de plus. On découpe
l'image ou la vidéo en petits morceaux. On remplace chaque morceau par une
liste de nombres, dans laquelle on inscrit aussi sa position. On empile des
blocs, où les jetons discutent entre eux puis réfléchissent chacun de leur
côté. Et l'on récupère à la sortie une file de jetons de la même taille
qu'à l'entrée, mais dont chacun a été enrichi par tous les autres.
\end{aretenir}

\vspace{0.5em}

Nous avons maintenant tout ce qu'il faut. Nous savons ce qu'est une image
pour une machine, ce que veut dire apprendre, comment jouer au jeu des trous
sans étiquettes, et de quoi est fait l'œil qui regarde. Ces quatre chapitres
formaient les fondations, et elles sont posées.

Une seule question reste en suspens, celle que nous avions laissée à la fin
du chapitre~\ref{chap:auto-supervision} : quand la machine devine ce qui est
caché, doit-elle deviner les pixels, ou l'idée ? La réponse à cette question
porte un nom, JEPA, et c'est par elle que s'ouvre la partie~II.

\phantomsection
\section*{Exercices}
\addcontentsline{toc}{section}{Exercices}

\begin{exercice}[Compter des jetons]
Vous donnez au modèle une image de $768 \times 768$ pixels, avec des patchs
de $16$. Combien de jetons obtenez-vous ? Et un clip de $32$ images en
$256 \times 256$, avec des tubelets de $2$ ?
\end{exercice}

\begin{solution}
Pour l'image : $768 / 16 = 48$ sur chaque côté, donc
$48 \times 48 = 2\,304$ jetons.

Pour le clip : $32 / 2 = 16$ tranches de temps, et $256/16 = 16$ patchs sur
chaque côté, donc $16 \times 16 \times 16 = 4\,096$ jetons. Remarquez que ce
clip coûte $(4\,096 / 2\,048)^2 = 4$ fois plus cher qu'un clip de la phase
principale, alors qu'il ne contient que deux fois plus d'images.
\end{solution}

\begin{exercice}[Choisir une résolution sous contrainte]
Votre serveur ne peut pas dépasser $5\,000$ jetons par clip. Vous voulez des
clips de $16$ images. Quelle est la plus grande résolution carrée, multiple
de $16$, que vous pouvez vous permettre ?
\end{exercice}

\begin{solution}
Avec $16$ images et des tubelets de $2$, il y a $8$ tranches de temps. Le
budget par tranche vaut donc $5\,000 / 8 = 625$ patchs. Comme l'image est
carrée, il faut au plus $\sqrt{625} = 25$ patchs de côté, ce qui donne
$25 \times 16 = 400$ pixels. Et $400$ est bien un multiple de $16$ :
la réponse est donc $400 \times 400$, qui consomme
$8 \times 25 \times 25 = 5\,000$ jetons, exactement le budget.

Un mot de prudence, quand même. Tomber pile sur la limite ne laisse aucune
marge, et l'article, lui, a retenu $384 \times 384$ pour ses vidéos, ce qui
donne $8 \times 24 \times 24 = 4\,608$ jetons. Choisir la résolution juste
en dessous de son plafond est une habitude qui évite bien des mauvaises
surprises le jour où le serveur est chargé.
\end{solution}

\begin{exercice}[Répartir entre les têtes]
Un modèle a des jetons de $1\,408$ nombres et $22$ têtes d'attention.
Combien de nombres chaque tête reçoit-elle ? Que se passerait-il avec $20$
têtes ?
\end{exercice}

\begin{solution}
Avec $22$ têtes : $1\,408 / 22 = 64$ nombres chacune. Le compte tombe juste.

Avec $20$ têtes : $1\,408 / 20 = 70{,}4$, ce qui n'est pas un nombre entier.
On ne peut pas donner $70{,}4$ nombres à une tête. La bibliothèque refusera
de construire le modèle. C'est pour cela que le nombre de têtes doit toujours
diviser exactement la taille du jeton.
\end{solution}

\begin{exercice}[Une softmax à la main]
Trois scores valent $1$, $1$ et $3$. Calculez les trois parts. Que
remarquez-vous sur les deux premières ?
\end{exercice}

\begin{solution}
Les exponentielles valent $e^1 \approx 2{,}718$, $e^1 \approx 2{,}718$ et
$e^3 \approx 20{,}086$. Leur somme vaut $25{,}522$. Les parts valent donc

\[
  \frac{2{,}718}{25{,}522} \approx 0{,}106, \qquad
  \frac{2{,}718}{25{,}522} \approx 0{,}106, \qquad
  \frac{20{,}086}{25{,}522} \approx 0{,}787 .
\]

Les deux premières sont égales, ce qui est normal : elles avaient le même
score. Et l'on voit l'effet de l'exponentielle : le score $3$ n'est que trois
fois plus grand que le score $1$, mais il récolte sept fois et demie plus de
part, puisque $0{,}787 / 0{,}106 \approx 7{,}4$. La softmax écarte bel et
bien les grands des petits.
\end{solution}

\begin{exercice}[Le mélange des positions]
Un collègue vous dit : «~j'ai mélangé l'ordre de mes jetons avant de les
donner au modèle, mais comme ils contiennent tous la même information
qu'avant, cela ne change rien~». A-t-il raison ?
\end{exercice}

\begin{solution}
Il a raison sur un point et tort sur l'essentiel. L'attention traite bien la
file comme un sac : à elle seule, elle rendrait le même résultat.

Mais les jetons ne contiennent pas que leur contenu : l'encodage de position
y a inscrit l'endroit d'où ils viennent. En mélangeant la file, il a donc
associé chaque morceau à une mauvaise adresse. Pour le modèle, c'est comme
si l'on avait recollé les pièces du puzzle au hasard : les pièces sont
justes, le dessin ne l'est plus. Le résultat sera faux, et sans le moindre
message d'erreur.
\end{solution}

\begin{exercice}[Le prix de la finesse]
Un ingénieur propose de passer les patchs de $16$ à $8$ pixels, «~pour voir
plus de détails~», en gardant des images de $256 \times 256$. Par combien le
travail de l'attention est-il multiplié ?
\end{exercice}

\begin{solution}
Avec des patchs de $16$, on avait $16 \times 16 = 256$ jetons. Avec des
patchs de $8$, on obtient $256/8 = 32$ patchs de côté, donc
$32 \times 32 = 1\,024$ jetons : quatre fois plus.

Comme le coût grimpe au carré, le travail est multiplié par
$4 \times 4 = 16$. Sa proposition améliorera sans doute la finesse, mais elle
coûtera seize fois plus cher. C'est un choix parfaitement défendable — à
condition de le faire les yeux ouverts, et de vérifier que le gain de qualité
vaut la dépense.
\end{solution}
